% ===========================================================================
\chapter{Introduction}
\label{chap:introduction}
% ===========================================================================

\section{Contexte}

Le rythme de livraison des logiciels s'est considérablement accéléré. Les
organisations informatiques sont soumises à des cycles de publication toujours
plus courts, tout en devant satisfaire des exigences non fonctionnelles
croissantes, en particulier en matière de sécurité et de conformité. L'adoption
de DevOps a constitué la réponse dominante à cette pression, en automatisant la
chaîne de construction et de déploiement~\cite{prates2025}.

Cette accélération a toutefois créé un déséquilibre. Lorsque la sécurité reste
une activité de fin de cycle, elle devient soit un goulot d'étranglement, soit
une étape escamotée. DevSecOps désigne l'intégration des pratiques de sécurité
dans DevOps~\cite{myrbakken2017}, selon le principe du \emph{shift security to
left} : faire intervenir la sécurité dès les premières étapes du cycle de
développement plutôt qu'au moment de la mise en production.

Le présent projet s'inscrit dans ce cadre. Il ne s'agit pas de développer une
application riche en fonctionnalités, puis de l'auditer. Il s'agit de concevoir
une chaîne de production logicielle dans laquelle les contrôles de sécurité sont
automatisés, outillés, et dont l'efficacité est démontrable.

\section{Problématique}

Le sujet du module formule une exigence précise quant à la nature de la
démonstration attendue : il ne suffit pas de présenter une liste d'outils, mais
d'établir pour chaque mécanisme de sécurité la chaîne complète

\begin{center}
\textbf{Pratique de sécurité $\rightarrow$ Outil $\rightarrow$ Intégration dans
la pipeline $\rightarrow$ Détection $\rightarrow$ Résultat $\rightarrow$ Action
corrective}
\end{center}

Cette exigence soulève une difficulté méthodologique qui a structuré l'ensemble
du projet et qu'il convient d'énoncer d'emblée.

\begin{constat}{une application sécurisée ne démontre rien}
Si le code produit est exempt de défaut dès le premier commit, chaque analyseur
de la pipeline retourne un rapport vide. Or un rapport vide ne démontre les
étapes \emph{Détection}, \emph{Résultat} et \emph{Action corrective} pour
aucune pratique.

La capacité à démontrer l'efficacité d'un contrôle suppose donc de disposer de
détections réelles. Nous avons choisi de les produire de manière volontaire et
contrôlée, selon la démarche décrite au chapitre~\ref{chap:vulnerabilites}.
\end{constat}

Une seconde difficulté concerne la nature même des pratiques à mettre en œuvre.
Comme le montre le chapitre~\ref{chap:fondements}, plusieurs des pratiques de
niveau Tier~1 ne sont pas outillables : elles relèvent du processus et de
l'organisation. Leur mise en œuvre ne produit pas de rapport d'outil mais des
artefacts documentaires et des configurations de plate-forme. Les traiter comme
des étapes de pipeline aurait conduit à un inventaire incomplet.

\section{Objectifs}

Le projet poursuit les objectifs suivants.

\begin{enumerate}
  \item \textbf{Identifier et mettre en œuvre les onze pratiques DevSecOps de
  niveau Tier~1} telles que définies par Prates et Pereira~\cite{prates2025},
  en automatisant leur exécution autant que possible.

  \item \textbf{Mettre en œuvre un sous-ensemble justifié de pratiques de
  niveau Tier~2}, en précisant pour chacune son objectif, l'outil retenu, son
  intégration dans la pipeline, les résultats obtenus et sa contribution à la
  sécurité.

  \item \textbf{Développer une application} comportant un front-end web, une
  application mobile, un ensemble d'API REST et un back-end métier, dont la
  sécurité est intégrée au cycle de développement.

  \item \textbf{Prendre en compte les catégories de cyber-risques} issues des
  référentiels OWASP : Top~10:2025 pour le web~\cite{owasptop10}, API
  Security Top~10~\cite{owaspapi}, Mobile Top~10~\cite{owaspmobile} et
  Top~10 for LLM Applications~\cite{owaspllm}.

  \item \textbf{Augmenter la pipeline par l'intelligence artificielle}, tant
  dans l'application que dans la chaîne d'intégration continue elle-même.

  \item \textbf{Démontrer l'efficacité de chaque contrôle} selon la chaîne de
  démonstration exigée, preuves à l'appui.

  \item \textbf{Analyser l'impact des contrôles de sécurité sur la performance
  de la pipeline}, et documenter les compromis retenus.
\end{enumerate}

\section{Périmètre et stratégie de réalisation}

\subsection{Priorité accordée à la profondeur plutôt qu'à l'étendue}

Le sujet précise que la priorité d'évaluation porte sur la qualité de la
démarche de cybersécurité, l'automatisation des contrôles et la capacité à
démontrer l'efficacité de la pipeline, plutôt que sur le nombre de
fonctionnalités développées.

Cette indication a été prise au sérieux. Le périmètre fonctionnel de
l'application est volontairement restreint, et chaque fonctionnalité retenue
l'a été parce qu'elle crée une surface de sécurité que nous devons ensuite
défendre, outiller et démontrer. Les fonctionnalités n'apportant aucune
catégorie de risque nouvelle ont été écartées, notamment le traitement réel des
paiements, l'envoi de courriels et l'authentification fédérée.

\subsection{Réalisation séquentielle du web puis du mobile}

L'application mobile, exigée par le sujet, a été délibérément reportée à une
seconde phase, postérieure à l'achèvement complet de l'application web et de sa
pipeline.

\begin{decision}{report de l'application mobile}
Une application web complète dotée d'une pipeline entièrement outillée et
documentée a davantage de valeur, au regard des critères d'évaluation énoncés,
que deux clients partiellement achevés partageant une pipeline superficielle.

Ce report n'est pas un abandon : l'architecture a été conçue dès l'origine
autour d'une API indifférente au client, de sorte que l'ajout du client mobile
ne nécessite aucune reprise côté serveur. Les travaux restants sont recensés
et planifiés.
\end{decision}

\subsection{Rédaction du rapport en parallèle de la réalisation}

Le présent rapport est rédigé au fur et à mesure de l'avancement, chaque phase
achevée étant documentée immédiatement. Ce choix répond à une contrainte
pratique : le raisonnement qui conduit à une décision technique est la première
chose qui s'efface. Les registres de décisions d'architecture
(\emph{Architecture Decision Records}) tenus tout au long du projet constituent
la matière première des sections de justification de ce document.

Cette démarche est par ailleurs cohérente avec les exigences de documentation
introduites par l'ASVS~5.0~\cite{owaspasvs5}, qui demande que certaines
décisions de sécurité soient explicitement consignées pour pouvoir être
vérifiées.

\section{Organisation du rapport}

Le chapitre~\ref{chap:fondements} présente les fondements de la démarche : le
passage de DevOps à DevSecOps, la taxonomie des pratiques retenue et les
référentiels de sécurité appliqués. Il identifie en particulier les onze
pratiques de niveau Tier~1.

Les chapitres~\ref{chap:archi-app} et~\ref{chap:archi-pipeline} décrivent
respectivement l'architecture de l'application et celle de la pipeline
DevSecOps.

Le chapitre~\ref{chap:phase0} documente la première phase de réalisation : la
mise en place de la gouvernance de sécurité du dépôt, antérieure à toute
écriture de code applicatif.

Les chapitres suivants documentent les phases ultérieures au fur et à mesure de
leur achèvement : modélisation des menaces, développement sécurisé, tests de
sécurité automatisés, étages de sécurité de la pipeline, conteneurisation,
supervision, augmentation par l'IA.

Le chapitre~\ref{chap:vulnerabilites} présente les démonstrations de
vulnérabilités volontaires et les corrections apportées. Le
chapitre~\ref{chap:performance} analyse l'impact des contrôles sur la
performance de la pipeline. Le chapitre~\ref{chap:limites} expose les limites du
travail et les perspectives.
